% Companion note 4 to "Relational Progression as an Operational Foundation for GR".
% Prepared with Claude (Anthropic) at the author's request.
\documentclass[11pt]{article}
\usepackage[margin=1in]{geometry}
\usepackage{amsmath,amssymb,amsthm}
\usepackage{booktabs,array}
\usepackage[shortlabels]{enumitem}
\usepackage{xcolor}
\usepackage[colorlinks=true,linkcolor=blue!50!black,citecolor=blue!50!black,urlcolor=blue!50!black]{hyperref}
\newcommand{\dd}{\mathrm{d}}
\newcommand{\st}[1]{\textsf{\small[#1]}}

\title{Testing the Inputs of Entanglement Equilibrium\\ on a 3D Lattice\\[4pt]
\large Companion note 4 to the consolidated draft}
\author{Scott A. Jordan}
\date{Draft, 23 November 2026}

\begin{document}
\sloppy
\maketitle

\begin{abstract}
Jacobson's entanglement-equilibrium argument derives the linearized Einstein equation, with $G=1/(4\hbar\eta)$, from three ingredients: (i) the ultraviolet entanglement entropy of a small ball changes as $\eta\,\delta A$ with the \emph{physical} area; (ii) the infrared entropy change of the matter fields is $\delta S_{\rm IR}=2\pi\int\! w\,\delta\langle T_{00}\rangle$ with the conformal ball weight $w=(R^2-r^2)/2R$; (iii) the total is stationary at fixed volume. On a fixed-background lattice for a free massless scalar in three space dimensions, (i) and (ii) can be computed and can fail; (iii) cannot be tested, because the geometry does not respond. We pre-registered both predictions and tested them. (i) holds: deforming the metric near the entangling sphere changes the entropy by $\kappa\,\delta A_{\rm phys}$ with the flat-space coefficient $\kappa=0.295$, to within $2\%$ for a smooth deformation and $5\%$ for a narrow one, independent of $R$; the coordinate area, which does not change, does not control the entropy. (ii) holds, with a sharpening: the exact lattice modular Hamiltonian of the ball matches $2\pi\int w\,T_{00}$ only when $T_{00}$ includes the bulk term of the \emph{improved} (conformal, $\xi=1/6$) stress tensor (its edge term is absent; note~5). A two-parameter fit gives the weight of the improvement term relative to the gradient term as $b/a=1.25,\,1.051,\,1.018,\,1.007$ for box sizes $N/R\approx2,4,6,8$, converging to the improved value $1$; the minimal stress tensor ($b/a=0$) is excluded. Both inputs of the argument are therefore derived numerically for this field. The equilibrium principle itself remains an assumption.
\end{abstract}

\section{What a fixed-background lattice can and cannot test}

In Jacobson's argument~\cite{Jacobson2016} a small geodesic ball of radius $R$ is varied away from the vacuum of flat space at fixed volume. The total entanglement entropy is split into
\begin{equation}
\delta S_{\rm tot}=\underbrace{\eta\,\delta A\big|_V}_{\text{UV, geometry}}+\underbrace{\delta S_{\rm IR}}_{\text{matter}},
\qquad \delta S_{\rm IR}=2\pi\!\int_{\rm ball}\! \frac{R^2-r^2}{2R}\,\delta\langle T_{00}\rangle\,\dd^3x .
\label{eq:split}
\end{equation}
Requiring $\delta S_{\rm tot}=0$ and using the geometric identity $\delta A|_V=-\tfrac{4\pi R^4}{15}G_{00}$ for a small ball gives $G_{00}=8\pi G\,T_{00}$ with $G=1/(4\hbar\eta)$; applying it in every frame gives the linearized Einstein equation. For a uniform deposit, $\delta S_{\rm IR}=\tfrac{8\pi^2R^4}{15}\delta T_{00}$, and the $R^4$ matches.

A lattice with a prescribed metric can compute the entropy of any state on any geometry, but the geometry never adjusts to the state. It follows that:
\begin{itemize}
\item \textbf{Testable (Test A).} Whether the entanglement entropy across a sphere responds to a change of the physical area with a definite coefficient $\eta$, and whether that coefficient is the flat-space area-law value. If the entropy tracked coordinate area, or changed with a different coefficient, there would be no well-defined $\eta$ to put in $G=1/(4\hbar\eta)$.
\item \textbf{Testable (Test B).} Whether the matter entropy change is the CHM form~\eqref{eq:split}, and with which local energy density. This is not automatic: the first law $\delta S=\delta\langle K\rangle$ holds for any state, but the claim that the modular Hamiltonian $K$ is the local operator $2\pi\int w\,T_{00}$ is a property of the vacuum that can fail on a lattice, or hold only with a particular $T_{00}$.
\item \textbf{Not testable.} The equilibrium principle $\delta S_{\rm tot}=0$, i.e.\ that the geometry \emph{chooses} $\delta A$ to cancel $\delta S_{\rm IR}$. On a fixed background $\delta A$ is put in by hand. Once (i) and (ii) hold, the statement ``the area change that restores equilibrium equals the GR value'' is the algebra above, not an independent check. Also not tested: the fixed-volume condition (Test A varies area and volume together), gravitons and other species, interacting fields, and the identification of the lattice spacing with a physical cutoff.
\end{itemize}

\section{Model}
A free massless scalar on $\mathbb{R}^3$, decomposed into angular sectors $(l,m)$ and discretised in the radial direction with spacing $1$ (Srednicki's radial lattice~\cite{Srednicki1993}), with Dirichlet wall at $r=N+1$. On a spatial metric $\dd\ell^2=\dd r^2+f(r)^2\dd\Omega^2$ the sector Hamiltonian is
\begin{equation}
H_l=\tfrac12\sum_{j=1}^{N}\Big[\pi_j^2+f_{j+\frac12}^2\Big(\frac{u_j}{f_j}-\frac{u_{j+1}}{f_{j+1}}\Big)^{\!2}+\frac{l(l+1)}{f_j^2}u_j^2\Big],
\end{equation}
with $u_j=f_j\,\psi_l(r_j)$; for $f=r$ this is Srednicki's chain. The ball is the first $n$ sites, $R=n+\tfrac12$. Ground-state correlators $X=\tfrac12K^{-1/2}$, $P=\tfrac12K^{1/2}$ give the reduced Gaussian state. Its entropy follows from the symplectic eigenvalues $\nu$ of $(X_AP_A)^{1/2}$, and its modular Hamiltonian $K_A=\tfrac12u^\top M_xu+\tfrac12\pi^\top M_p\pi$ from
\begin{equation}
M_x=X_A^{-1/2}Q\,\mathrm{diag}(\nu\varepsilon)\,Q^\top X_A^{-1/2},\quad
M_p=X_A^{1/2}Q\,\mathrm{diag}(\varepsilon/\nu)\,Q^\top X_A^{1/2},\quad \varepsilon=\ln\frac{2\nu+1}{2\nu-1},
\end{equation}
where $X_A^{1/2}P_AX_A^{1/2}=Q\,\mathrm{diag}(\nu^2)Q^\top$. Because $\nu-\tfrac12$ reaches $10^{-92}$, Test B is computed with 100-digit arithmetic (\texttt{mpmath}).

\section{Test A: response to physical area}
\paragraph{Pre-registered prediction.} For a deformation $f(r)=r\,[1+\epsilon\, e^{-(r-R)^2/w^2}]$ that leaves the coordinate sphere and lattice unchanged but changes the physical area of the entangling sphere, $\delta S=\kappa\,[f(R)^2-R^2]$ (area in units of $4\pi a^2$, $\kappa=0.295$ from the flat-space fit). \emph{Contradiction:} $\delta S\approx0$ (entropy tracks coordinate area) or a ratio $\delta S/\kappa\delta A$ that differs from $1$ by more than the expected curvature corrections, or drifts with $R$.

\paragraph{Result} ($N=60$, angular sum to $l=3000$):
\begin{center}\small
\begin{tabular}{ccccc}\toprule
$R$ & $\epsilon=+0.02,\ w=3$ & $\epsilon=-0.02,\ w=3$ & $\epsilon=+0.05,\ w=3$ & $\epsilon=+0.02,\ w=6$\\\midrule
10.5 & 1.050 & 1.049 & 1.051 & 1.015\\
15.5 & 1.053 & 1.051 & 1.054 & 1.018\\
20.5 & 1.053 & 1.052 & 1.054 & 1.019\\
25.5 & 1.054 & 1.052 & 1.055 & 1.019\\\bottomrule
\end{tabular}
\end{center}
Entries are $\delta S/(\kappa\,\delta A)$. Example: $R=25.5$, $\epsilon=0.02$, $w=3$: $\delta S=+8.165$, prediction $+7.750$. The response is linear in $\epsilon$ (sign-symmetric, unchanged at $\epsilon=0.05$), independent of $R$, and approaches $1$ as the deformation is made smoother ($5\%$ at $w=3$, $2\%$ at $w=6$). The residual is what a curvature or extrinsic-curvature correction from the localized bump would give; it is not a drift of the area coefficient. \textbf{Supported.}

\paragraph{Limits.} Angular cutoff: at $R=20.5$ the ratio at $l_{\max}=750,1500,3000,6000$ is $1.0478,1.0522,1.0534,1.0538$ ($w=3$) and $1.0132,1.0176,1.0189,1.0192$ ($w=6$), so the remaining tail beyond $l=3000$ is below $0.1\%$. Box size: doubling the box to $N=120$ at $R=20.5$ changes the ratios by $2\times10^{-4}$ ($1.0539$ and $1.0194$ at $l_{\max}=6000$), so the area response is a short-distance effect insensitive to the wall. The deformation changes volume as well as area, so the fixed-volume condition of~\eqref{eq:split} is not probed.

\section{Test B: the CHM modular Hamiltonian, and which stress tensor}
\paragraph{Probe.} A classical-noise deposit $X_A\to X_A+\epsilon\,vv^\top$ with $v_j=r_j\psi(r_j)$ and $\psi$ a Gaussian bump of centre $r_0$ and width $\sigma$, strictly inside the ball. Then $\delta S/\epsilon\to Q_{\rm exact}=\tfrac12v^\top M_xv$ (checked against the exact entropy: $772.96387$ vs $772.9677$ for one bump). The two candidate local forms are
\begin{align}
Q_{\min}&=2\pi\sum_j w\big(r_{j+\frac12}\big)\,\tfrac12(j+\tfrac12)^2\Big(\tfrac{v_j}{j}-\tfrac{v_{j+1}}{j+1}\Big)^{\!2}+2\pi\sum_j w(r_j)\frac{l(l+1)}{2j^2}v_j^2,\\
Q_{\rm imp}&=Q_{\min}+\frac{\pi}{R}\sum_j v_j^2 ,
\end{align}
where the extra term is $2\pi\int w\,(-\xi\nabla^2\psi^2)$ with $\xi=1/6$ and $\nabla^2w=-3/R$; the boundary term vanishes because the bump vanishes at $R$. A momentum deposit, $\tfrac12p^\top M_pp$ against $2\pi\sum w\,p_j^2/2$, involves no improvement term and serves as a control.

\paragraph{Pre-registered prediction.} $Q_{\rm exact}=Q_{\rm imp}$; $Q_{\min}$ falls short by the predicted $\tfrac{\pi}{R}\sum v^2$. \emph{Contradiction:} $Q_{\rm exact}$ matches neither within the finite-box error measured by the control, or matches $Q_{\min}$.

\paragraph{Result} ($N=96$, $n=24$, $R=24.5$):
\begin{center}\small
\begin{tabular}{c|ccc|ccc|ccc}\toprule
& \multicolumn{3}{c|}{$l=0$} & \multicolumn{3}{c|}{$l=1$} & \multicolumn{3}{c}{$l=2$}\\
$(r_0,\sigma)$ & $\frac{Q_{\rm ex}}{Q_{\min}}$ & $\frac{Q_{\rm ex}}{Q_{\rm imp}}$ & $\frac{Q_{\rm imp}}{Q_{\min}}$ & $\frac{Q_{\rm ex}}{Q_{\min}}$ & $\frac{Q_{\rm ex}}{Q_{\rm imp}}$ & $\frac{Q_{\rm imp}}{Q_{\min}}$ & $\frac{Q_{\rm ex}}{Q_{\min}}$ & $\frac{Q_{\rm ex}}{Q_{\rm imp}}$ & $\frac{Q_{\rm imp}}{Q_{\min}}$\\\midrule
(0,4) & 1.095 & 1.054 & 1.038 & 1.049 & 1.037 & 1.011 & 1.014 & 1.009 & 1.005\\
(0,7) & 1.196 & 1.052 & 1.137 & 1.061 & 1.026 & 1.035 & 1.022 & 1.008 & 1.014\\
(8,3) & 1.120 & 1.049 & 1.068 & 1.062 & 1.017 & 1.044 & 1.031 & 1.005 & 1.026\\
(12,3) & 1.133 & 1.042 & 1.088 & 1.083 & 1.012 & 1.069 & 1.051 & 1.002 & 1.049\\
(12,5) & 1.348 & 1.042 & 1.294 & 1.168 & 1.007 & 1.160 & 1.084 & 1.000 & 1.084\\
(16,2) & 1.087 & 1.034 & 1.052 & 1.061 & 1.012 & 1.048 & 1.046 & 1.003 & 1.043\\\midrule
control $M_p$ & \multicolumn{3}{c|}{1.042--1.049} & \multicolumn{3}{c|}{1.008--1.015} & \multicolumn{3}{c}{1.001--1.005}\\\bottomrule
\end{tabular}
\end{center}
In all 18 cases $Q_{\rm exact}$ is closer to $Q_{\rm imp}$, and the deviation from $Q_{\min}$ follows the predicted $Q_{\rm imp}/Q_{\min}$ (e.g.\ $1.348$ observed vs $1.294$ predicted for the broad bump, where the improvement term is largest). The remaining excess $Q_{\rm ex}/Q_{\rm imp}-1$ is the same size as the momentum control, which has no improvement term; so it is a common finite-box factor, not a stress-tensor effect.

\paragraph{Separating the box factor.} To remove the common factor we fit $Q_{\rm exact}=a\,Q_{\min}+b\,(Q_{\rm imp}-Q_{\min})$ over the six bumps (this fit was added after the first run; the raw comparison above was the pre-registered one). The improved tensor predicts $b/a=1$, the minimal $b/a=0$. For $l=0$, $R=24.5$:
\begin{center}\small
\begin{tabular}{cccccc}\toprule
$N$ & $N/R$ & $a$ & $b/a$ & max fit residual & momentum control\\\midrule
48 & 2.0 & 1.153 & 1.249 & $8\times10^{-2}$ & 1.20--1.24\\
96 & 3.9 & 1.033 & 1.051 & $1.8\times10^{-2}$ & 1.042--1.049\\
144 & 5.9 & 1.015 & 1.018 & $7.6\times10^{-3}$ & 1.018--1.021\\
192 & 7.8 & 1.009 & 1.007 & $3.8\times10^{-3}$ & 1.010--1.012\\\bottomrule
\end{tabular}
\end{center}
Both $a$ and $b/a$ converge to $1$ as the wall recedes, with $a-1$ tracking the momentum control. For $l=1$ and $l=2$ at $N=144$, $b/a=0.962$ and $1.030$; there the improvement term is a smaller share of the total ($1$--$16\%$), so $b$ is less tightly fixed, but both are within $4\%$ of $1$ and far from $0$. At $R=12.5$, $N=96$ and $144$ give $b/a=0.988$ and $0.973$ with bumps of width $1$--$2$ sites, where the improvement term is only $1$--$5\%$ of the total and $b$ is poorly constrained. \textbf{Supported, with the improved stress tensor; minimal excluded.}

\paragraph{Limits.} Precision: 100 digits; the exact entropy agrees with the first-law value to $5\times10^{-6}$. Finite box: a $1\%$ residual at $N/R\approx8$, falling roughly as $(R/N)^2$. Lattice: bumps of width $\ge2$ sites; the comparison near the origin and near $R$ (within $\sim2$ sites) is not tested. Only the sectors $l\le2$ are checked; since the vacuum and the CHM operator both decompose into $(l,m)$ sectors, each sector is an independent test.

\section{What this means for the program}
\begin{enumerate}
\item \textbf{Ingredient (i) is derived numerically} for a free scalar: its entanglement entropy is a function of physical area with the flat coefficient, so $\eta=\kappa/a^2$ is well defined and the earlier identification $G=1/(4\hbar\eta)$, $a\approx0.31\sqrt{N_{\rm sp}}\,\ell_P$, is not undermined by coordinate dependence. \st{Derived (numerical)}
\item \textbf{Ingredient (ii) is derived numerically}, and sharpened: the matter side of entanglement equilibrium is weighted by the \emph{conformal} stress tensor, even though the lattice field is minimally coupled. The extra $\tfrac{\pi}{R}\int\psi^2$ is an expectation value of $\phi^2$, a scalar operator of dimension $1$. In Jacobson's argument, departures from the conformal form are written as a term $\delta X$ built from such operators~\cite{Jacobson2016}, and Casini, Galante and Myers found that relevant operators of low dimension threaten that step~\cite{CasiniGalanteMyers2016}. Here the lattice shows that, for the free scalar, the whole $\phi^2$ term is accounted for by calling $T_{00}$ the improved ($\xi=1/6$) energy density, with no further correction visible at the $1\%$ level. Integrated over all space the improvement is a total derivative, so total mass is unaffected. Locally, though, the source that entanglement equilibrium would hand to the geometry is $T^{\rm imp}_{ab}$, not the minimally coupled $T_{ab}$ that ordinary GR would use for this field. For a smooth deposit the extra term scales as $R^2$, so at small $R$ it dominates the $R^4$ energy term rather than being a small correction. \emph{Correction (note 5):} a follow-up test with bumps that reach the sphere shows that the improved tensor's edge term $-\tfrac{\pi}{3}\oint\phi^2$ is \emph{absent}, so the tests here establish the bulk improvement term only; see companion note~5 for the consequences. \st{Derived (numerical)}
\item \textbf{Ingredient (iii) remains assumed.} No fixed-background computation can test it. What would test it is a model in which the geometry is itself a quantum variable, so that $\delta A$ is determined by the state rather than prescribed. \st{Assumed}
\end{enumerate}
The status of $G$ from the environment is therefore unchanged in kind: it is ``Open (narrowed)'', now narrowed to the equilibrium principle alone, with both of its inputs computed.

\section*{Code}
\texttt{eq\_testA.py} (Test A; argument $N$), \texttt{eq\_testA\_fs.py} (angular-cutoff and box-size check; arguments $N$, $n$), \texttt{eq\_testB.py} (Test B; arguments: digits, $N$, $n$, comma-separated $l$ list). Runtimes: Test A about 2 minutes per $N=60$ table; Test B from 40\,s ($N=96$, one sector) to 5 minutes ($N=192$).

\begin{thebibliography}{9}
\bibitem{Jacobson2016} T.~Jacobson, ``Entanglement equilibrium and the Einstein equation,'' \emph{Phys. Rev. Lett.} \textbf{116}, 201101 (2016).
\bibitem{CasiniGalanteMyers2016} H.~Casini, D.~A.~Galante, and R.~C.~Myers, ``Comments on Jacobson's `Entanglement equilibrium and the Einstein equation','' \emph{JHEP} \textbf{03}, 194 (2016), arXiv:1601.00528.
\bibitem{Srednicki1993} M.~Srednicki, ``Entropy and area,'' \emph{Phys. Rev. Lett.} \textbf{71}, 666 (1993).
\end{thebibliography}
\end{document}
